% =====================================================================
% RETRIEVAL / POSE ENCODER SECTION
% Bashar Kabbarah — draft for OpenTouch main submission
%
% TODO markers flag every value I could not confirm from HANDOFF.md.
% Do not submit with any \TODO surviving.
%
% VERIFICATION PASS 2026-08-19: every value below re-derived from result
% JSONs / code / run params.txt, never from HANDOFF.md. Source citations
% sit in % comments beside each value. ZERO \TODO markers remain: the four
% flagged claims were resolved (MANO -> hand keypoints; the ReLU is now
% stated for the GRU arm and the no-nonlinearity fact scoped to the
% baseline; 10.25 softened to "roughly 10" pending an artifact; the
% harness-records-split sentence was made true by patching eval.py).
% =====================================================================

\newcommand{\TODO}[1]{\textcolor{red}{[TODO: #1]}}

% ---------------------------------------------------------------------
% METHODS  (main body — target ~1/3 column)
% ---------------------------------------------------------------------

\subsection{Pose encoder}
\label{sec:pose-encoder}

% T=20 per every retrieval run's params.txt (sequence_length: 20) and
% scripts/train.sh:17; eval.py:110 default is also 20. The draft's T=36 was
% the regression/probe window, never the retrieval window.
% "hand keypoints", deliberately NOT "MANO": OpenTouch pose is glove-derived
% (21-keypoint MediaPipe-style layout, wrist at index 0 -- see
% src/opentouch/pose_encoder.py:6-9); no MANO model is involved anywhere in
% the pose pipeline, and a reviewer who knows MANO's joint set would call
% the misattribution out.
Each pose sequence is a window of $T=20$ frames of 3D hand keypoints,
$21$ joints $\times\,3$ coordinates, giving a $63$-dimensional vector per
frame. The pose branch maps this window to a single embedding in the shared
tactile--pose space.

Prior versions of this encoder average-pooled the per-frame features and
projected the pooled $128$-dimensional summary directly to the shared space,
% src/opentouch/pose_encoder.py:76,115 -- Linear(128, 64), no nonlinearity
% on the mean path; discarding temporal order is the point of the baseline
discarding the temporal ordering of the window. We instead encode
the sequence with a bidirectional GRU (hidden size $120$, $2$ layers),
% src/opentouch/pose_encoder.py:56 -- nn.GRU(128, 120, num_layers=2, bidirectional=True)
take the concatenation of the final forward and final backward hidden states
% src/opentouch/pose_encoder.py:105
as a $240$-dimensional summary, and apply a ReLU followed by a single linear
% ReLU: src/opentouch/pose_encoder.py:108 (GRU branch only -- the baseline
% has none, see pose_encoder.py:110-115 and the baseline-fidelity paragraph)
projection to the $64$-dimensional shared space.
% src/opentouch/pose_encoder.py:76 -- Linear(240, 64); embed_dim=64 per
% src/opentouch/model_configs/OpenTouch-DINOv3-B16-Retrieval.json

Tactile and pose branches are trained jointly with a symmetric InfoNCE
% symmetric: src/opentouch/loss.py:106-109 (mean of CE over both directions)
objective over in-batch negatives, at a learnable temperature initialised to
$0.07$,
% src/opentouch/model.py:81-84 -- logit_scale nn.Parameter, init log(1/0.07)
using AdamW with learning rate $10^{-4}$, batch size $256$,
for $300$ epochs. Results are reported over three encoder seeds
($42$, $0$, $1$).
% src/opentouch_train/main.py:235-242 (AdamW; betas 0.9/0.999 and eps 1e-8
% from params.py:7; wd 0.01 excluded for bias/BN/logit_scale). Schedule:
% cosine with 5% linear warmup (main.py:286-291; every run's params.txt:
% lr 0.0001, batch_size 256, epochs 300, lr_scheduler cosine, warmup 0.05).
% Seed runs identified EMPIRICALLY by re-evaluation 2026-08-19:
%   GRU:     logs/2026_06_22-21_01_54 (s42), 2026_07_06-22_35_40 (s0),
%            2026_07_06-22_36_12 (s1)
%   AvgPool: logs/clip_avgpool_s{42,0,1}
% (the 2026_07_06-15_2x pair scores ~72 on the clip gallery -- trained
% pre-028eba5 when --seed also reseeded the SPLIT, so different partition;
% quarantined on the cluster. Full mapping, evidence and quarantine record:
% results/RETRIEVAL_SEED_PROVENANCE.md)

\paragraph{Retrieval protocol.}
We report mean average precision as the mean reciprocal rank of the correct
counterpart, scored against the entire evaluation split as the gallery.
% src/opentouch/metrics.py:75-77 -- single positive (diagonal), rank =
% #(sims >= correct sim) over the full split, mAP = mean(1/rank). With one
% positive per query, AP == reciprocal rank, so the sentence is exact.
This metric is a function of gallery size, so we report the number of gallery
windows alongside every mAP. We evaluate under two splits. The
\emph{clip-disjoint} split holds out whole clips but allows the same
participant to appear in train and evaluation. The
\emph{participant-disjoint} split holds out whole scenes, so no evaluation
participant is seen during training.

% ---------------------------------------------------------------------
% APPENDIX
% ---------------------------------------------------------------------

\subsection{Pose encoder: retrieval results}
\label{app:retrieval}

\begin{table}[h]
\centering
\caption{Tactile--pose retrieval, clip-disjoint split, test. Mean $\pm$ s.d.
over three seeds ($42$, $0$, $1$); gallery $N=1399$ windows over $296$ clips.}
% gallery: results_bootstrap_clip_gru_test.json tactile_to_pose.n_test_samples
% = 1399, .n_clusters = 296 (same in all three avgpool seed bootstrap JSONs)
\label{tab:retrieval-clip}
\begin{tabular}{lccc}
\toprule
Pose encoder & T$\rightarrow$P mAP & P$\rightarrow$T mAP & Ratio \\
\midrule
% point estimates, 3 seeds: results/results_retrieval_clip_avgpool_s{42,0,1}_test.json
% T2P 16.2772/16.9798/16.1288, P2T 15.6856/16.2723/16.3113 (mean +- population sd)
Average pooling & $16.46 \pm 0.37$ & $16.09 \pm 0.29$ & --- \\
% point estimates, 3 seeds: results/results_retrieval_clip_gru_s{42,0,1}_test.json
% T2P 45.4634/46.2843/46.7619, P2T 44.3875/46.7151/45.7892
% ratios of means: 2.8046 (T2P), 2.8360 (P2T)
Sequence (biGRU) & $46.17 \pm 0.54$ & $45.63 \pm 0.96$ & $2.80\times$ / $2.84\times$ \\
\bottomrule
\end{tabular}
\end{table}

\begin{table}[h]
\centering
\caption{Tactile--pose retrieval, participant-disjoint split, test.
Point estimates with clip-clustered bootstrap 95\% intervals
($1000$ draws, fixed gallery, queries resampled); gallery $N=1411$ windows
over $257$ clips drawn from the held-out scenes (the bootstrap cluster unit
is the clip under both splits).}
% gallery: results_bootstrap_scene_gru_test.json tactile_to_pose
% .n_test_samples = 1411, .n_clusters = 257; protocol fields: .n_bootstrap =
% 1000, .cluster_unit = "clip", .gallery = "fixed (queries resampled only)"
\label{tab:retrieval-scene}
\begin{tabular}{lcc}
\toprule
Pose encoder & T$\rightarrow$P mAP & 95\% CI \\
\midrule
% point estimate: results/results_retrieval_scene_avgpool_norelu_test.json
% .tactile_to_pose_mAP = 0.064422 (draft's 6.46 was the bootstrap MEAN, which
% contradicted this caption's "point estimates"); CI from
% results/results_bootstrap_scene_avgpool_test.json
Average pooling & $6.44$ & $[5.45,\ 7.59]$ \\
% point estimate: results/results_retrieval_scene_p2t_scene_gru_test.json
% .tactile_to_pose_mAP = 0.283080; CI from results/results_bootstrap_scene_gru_test.json
Sequence (biGRU) & $28.31$ & $[25.59,\ 31.01]$ \\
\bottomrule
\end{tabular}
\end{table}

Replacing average pooling with a sequence encoder improves tactile--pose
retrieval by $2.80\times$ (T$\rightarrow$P) and $2.84\times$
(P$\rightarrow$T) on the clip-disjoint split. The margin widens to
% 0.283080 / 0.064422 = 4.3941 from the unrounded point estimates. The
% draft's 4.40 arises only by dividing already-rounded values (28.31/6.44)
% and is not reproducible from any consistent unrounded pairing.
$4.39\times$ under participant-disjoint evaluation, where the bootstrap
intervals for the two encoders do not approach overlap.

Both absolute figures fall sharply when participants are held out
% 46.17 (3-seed mean; the draft's 45.5 was seed 42 alone) -> 28.31 point;
% 16.46 (3-seed mean) -> 6.44 point
($46.2 \rightarrow 28.3$ for the sequence encoder, $16.5 \rightarrow 6.4$
for average pooling), indicating that a substantial part of clip-disjoint
retrieval performance reflects participant-specific regularities rather than
transferable tactile--pose structure. The architectural conclusion is
unaffected by this and in fact strengthens: temporal structure in the pose
window matters more, not less, when the model cannot rely on having seen the
participant.

\paragraph{Baseline fidelity.}
The average-pooling baseline reported here is the upstream architecture
exactly: a linear projection applied directly to the $128$-dimensional pooled
% src/opentouch/pose_encoder.py:76,115 -- Linear(128, 64), no ReLU on this path
vector. Two discrepancies were found and corrected before these numbers were
produced. First, an earlier implementation zero-padded the pooled vector to
$240$ dimensions so that both pooling modes could share a projection shape,
% commits f147054 (padded version) -> f71a25e (fix); HANDOFF.md:305-311
which prevented upstream checkpoints from loading. Second, a ReLU that was
introduced alongside the sequence encoder had been inherited by the pooling
path, where it does not belong; with that ReLU present, the upstream
% commit 1ab43d4; the no-ReLU requirement is pinned by
% tests/test_pose_encoder_temporal_mode.py
checkpoint scores roughly $10$ rather than $16.76$.
% "roughly 10" is a deliberate softening: the precise 10.25 is documented
% (HANDOFF.md:315, VALIDATION.md:681, and pose_encoder.py:29's own comment
% says "roughly 10") but has no committed eval artifact, and regenerating it
% would require the pre-fix encoder (commit f147054) plus the orphan
% upstream checkpoint, whose training provenance was never identified
% (results/RETRIEVAL_SEED_PROVENANCE.md). The architectural fact the number
% illustrates -- mean mode must not have the ReLU -- is test-pinned either
% way. 16.76: results/results_retrieval_clip_avgpool_test.json
% .tactile_to_pose_mAP = 0.167565 (the orphan run; the seeded baselines in
% Table 1 supersede it for all comparisons).
Both fixes make the baseline
stronger, and all comparisons above use the corrected baseline.

\paragraph{Reproducibility note.}
Retrieval mAP is computed against the full evaluation gallery, so the split
used to construct that gallery must match the split the checkpoint was
trained under. Evaluating a participant-disjoint checkpoint against a
clip-disjoint gallery inflates T$\rightarrow$P mAP to $72.09$,
% regenerated deliberately 2026-08-19, same bootstrap protocol and seed:
% results/results_bootstrap_scene_gru_clipgallery_demo_test.json
% .tactile_to_pose.mean_mAP = 72.088 (vs correct-gallery 28.31)
because the gallery then contains the model's own training participants.
The evaluation harness records the gallery split alongside every result.
% This sentence was overbroad when drafted (only bootstrap_eval.py recorded
% split_group_by) and was made TRUE rather than weakened: as of the same
% commit as this text, opentouch_train/eval.py writes a _meta block
% (checkpoint, split, split_group_by + source, split seed) into every output
% JSON. The per-seed JSONs cited in Table 1 already carry it. Caveat that
% does not affect the sentence: results_retrieval_*.json files produced
% BEFORE this change record metrics only.
